\documentclass[11pt,a4paper]{ctexart}
\usepackage[margin=2.4cm]{geometry}
\usepackage{amsmath,amssymb}
\usepackage{fancyhdr}
\usepackage{enumitem}
\linespread{1.18}
\setlength{\headheight}{15pt}
\setlength{\parskip}{0pt}

\pagestyle{fancy}
\fancyhf{}
\fancyhead[L]{\small Team 2407093}
\fancyhead[R]{\small \thepage}
\renewcommand{\headrulewidth}{0.4pt}
\fancypagestyle{plain}{\fancyhf{}\renewcommand{\headrulewidth}{0pt}}

\ctexset{
  section/format = \Large\bfseries\raggedright,
  subsection/format = \large\bfseries\raggedright,
  subsubsection/format = \normalsize\bfseries\raggedright
}
\setcounter{secnumdepth}{3}

\begin{document}

%====================== 摘要页 ======================
\thispagestyle{empty}
\begin{center}
{\small 2024 MCM/ICM\quad 选题 A\quad 队伍控制号 2407093\quad 摘要页（Summary Sheet）}\\[18pt]
{\LARGE\bfseries 共存还是灭绝？七鳃鳗与环境的关系}\\[8pt]
{\large （Coexist or Extinct? Relationship between Lampreys and Environment）}\\[10pt]
\end{center}

\begin{center}
{\large\bfseries 摘\quad 要}
\end{center}

\vspace{6pt}
与大多数其他动物不同，七鳃鳗能够响应环境变化而改变自身的性别比例。本文建立了 BFM 模型，用以解释具有可变性别比例的七鳃鳗对环境变化所表现出的高度适应性。

在动手完成各项任务之前，我们做了充分的准备工作。我们把七鳃鳗种群划分为三类：幼体、成年雌性与成年雄性。我们首先引入逻辑斯谛模型（Logistic Model）来刻画七鳃鳗的增长过程；随后引入性别分化比例 $\alpha$ 与幼体七鳃鳗的发育时长 $\tau$，用以描述种群的性别结构与年龄结构。由此我们得到 BFM 模型，并对其进行了初步验证。我们发现任务 1 与任务 4 关系密切，为使任务完成得更为顺畅，我们按照 1、4、2、3 的顺序依次作答。

关于任务 1 与任务 4：我们首先分别将 BFM 模型中的 $\alpha$ 与 $\tau$ 设为常数，以及设为随时间变化的函数。随后引入洛特卡–沃尔泰拉模型（Lotka-Volterra Model）与尼科尔森–贝利模型（Nicholson-Bailey Model），用以刻画七鳃鳗与其天敌、宿主及寄生虫之间的关系。接着我们比较了两种情形下的仿真结果。我们发现，七鳃鳗种群的可变性别比例能够降低生态系统中其他物种的灭绝风险。通过计算辛普森多样性指数（Simpson's diversity index）与香农–维纳指数（Shannon-Weiner index）可知，七鳃鳗的可变性别比例能够提升更大生态系统的物种多样性。

关于任务 2：为讨论可变性别的优缺点，我们通过改变成年七鳃鳗的死亡率 $\mu$ 进行数值仿真，使七鳃鳗种群分别处于适宜环境与扰动环境之中。结果表明，适应性的性别比例变化可使七鳃鳗种群在两种环境中都更加稳定。随后我们定义繁殖成功率来刻画七鳃鳗的基因传递，发现改变性别比例的能力并不利于七鳃鳗种群自身的演化。我们引入反馈模型（Feedback Model）解释适应性性别比例提升种群稳定性的内在机制，并绘制反馈机制图与相轨迹图，从定性与定量两个层面对稳定性加以说明。

关于任务 3：为更好地评价生态系统稳定性，我们引入抵抗力、恢复力与生态系统综合稳定性三个概念，并结合我们的模型给出定量描述。同时，我们还将物种多样性作为衡量生态系统稳定性的指标。我们的结论是：可变性别能够提升生态系统稳定性。

最后，我们对模型进行了灵敏度分析。与此同时，我们针对七鳃鳗入侵的治理提出了若干建设性建议，并分析了模型的优缺点、讨论了模型的改进方向。

\vspace{6pt}
\noindent\textbf{关键词：}反馈，尼科尔森–贝利模型，生态系统稳定性，生物入侵

\newpage
\pagestyle{fancy}

%%====================== 正文 ======================
\section{引言}

\subsection{问题背景}

过去一年间，五大湖（Great Lakes）的海七鳃鳗控制计划取得了巨大成功。该计划保护了当地的水产养殖业，使其免受入侵性海七鳃鳗的危害。海七鳃鳗在生态系统中扮演着复杂的角色：它们既是渔业的威胁，又是食物链中的关键一环，甚至还是一道珍馐。它们的独特性或许正与其改变性别比例的能力相关，这正是我们接下来要探讨的内容。

\subsection{问题重述}

\begin{itemize}
\item 首先，我们需要建立一个数学模型，用以描述七鳃鳗单独存在时的增长与繁殖条件。
\item 我们需要扩展模型，以考虑七鳃鳗与其他物种的相互作用，进而回答七鳃鳗对更大生态系统的影响。
\item 此后，我们应当比较性别比例可变与性别比例恒定时七鳃鳗种群的生长状况，从而得到具有可变性别比例的七鳃鳗种群的优缺点。
\item 进一步，我们需要分析适应性的性别比例变化对生态系统稳定性的影响。
\item 最后，我们需要探究七鳃鳗能够为生态系统内其他物种提供怎样的生存优势。
\end{itemize}

\subsection{文献综述}

在研究的起点，我们梳理了三类与性别分化相关的方法：分子生物学方法、种群统计学方法与性选择博弈论方法。其中，分子生物学方法着眼于分子层面性别分化的内在机制；性选择博弈论方法[2]则着重描述物种在收益最大化驱动下所采取的性别分化策略。然而，这两类方法都不适合用来描述生态系统中与环境相互作用的七鳃鳗种群。

因此我们认为，最恰当的做法是把种群统计学与生态学理论结合起来，对七鳃鳗身上观察到的性别分化现象进行数学建模。

\subsection{本文工作}

本文的整体研究框架如图 2 所示，后文各节依次展开其中的各个环节。

\section{模型准备}

\subsection{假设与合理性说明}

\noindent\textbf{假设 1：}七鳃鳗的生命周期仅包含在幼体期与成体寄生期这两个具有持续时间的阶段；从卵中孵化、性别分化与交配等阶段均视为瞬时完成。

\noindent\textbf{理由：}根据文献[3]，与七鳃鳗的幼体阶段和成体阶段相比，其生命周期的其他阶段持续时间相对较短。为使模型更聚焦于主要内容，我们认为只有幼体阶段与成体阶段具有显著的时间跨度。

\noindent\textbf{假设 2：}幼体七鳃鳗视为无性别之分，而成体七鳃鳗分化为雌、雄两种性别。

\noindent\textbf{理由：}幼体七鳃鳗的生殖系统尚未完全分化，七鳃鳗幼体处于生长阶段，尚未性成熟，也不参与繁殖。因此幼体七鳃鳗的性别可以不予考虑。

\noindent\textbf{假设 3：}幼体阶段与成体阶段的能量来源不同；对处于同一空间内的全部幼体七鳃鳗而言，单位时间内可获得的能量总量为常数。

\noindent\textbf{理由：}根据文献[4]，幼体七鳃鳗在湖底或海底营滤食性生活，主要以水体中的有机碎屑为食；而成体则在水中营半寄生生活，主要从宿主的营养物质中获取能量。随河流水流携带的有机碎屑在水体中应当均匀分布，在时间与空间上不会发生显著变化。因此可以假定幼体七鳃鳗单位时间内获得的能量来源恒定。

\noindent\textbf{假设 4：}七鳃鳗的繁殖仅与雌性有关。

\noindent\textbf{理由：}根据文献[5]，考虑到七鳃鳗的群体交配行为中卵的数量显著少于精子的数量，可以认为最终受精卵的数量主要取决于雌性产卵的数量。

\subsection{符号说明}

本文所用的主要符号及其含义如下：$B$ 表示幼体七鳃鳗的种群数量；$F$ 表示成年雌性七鳃鳗的种群数量；$M$ 表示成年雄性七鳃鳗的种群数量；$r$ 表示幼体七鳃鳗的增长率；$K$ 表示幼体七鳃鳗的环境容纳量；$\mu$ 表示死亡率；$\tau$ 表示幼体期的持续时间；$\alpha$ 表示性别分化比例；$P$ 表示捕食者种群数量；$H$ 表示宿主种群数量；$r_H$ 表示宿主的自然增长率；$J$ 表示寄生虫种群数量。

\subsection{七鳃鳗的生活史}

七鳃鳗的生活史如图 3 所示。

\section{BFM 模型的建立}

\subsection{微分方程形式}

首先考虑理想情形，即七鳃鳗生活在空间、食物及其他资源均不受限制的环境中。设 $r'$ 表示七鳃鳗的增长率，则其种群数量 $N$ 应满足如下方程：
\[
\frac{dN}{dt}=r'N
\]
由于空间、温度与食物等资源的限制，增长率会随种群规模的增大而下降。因此我们采用逻辑斯谛种群增长模型来描述种群增长过程，上式应改写为：
\[
\frac{dN}{dt}=r'N\left(1-\frac{N}{K'}\right)
\]
根据假设 2 与假设 4，幼体七鳃鳗的出生率仅与雌性的数量有关；且幼体与成体七鳃鳗栖息于不同区域，二者的环境容纳量相互独立。然而，一部分幼体七鳃鳗会成长为成体，从而使幼体群体的数量减少。设 $\bar B$ 表示成长为成体七鳃鳗的数量，结合文献[5]的结论，我们采用如下微分方程描述幼体七鳃鳗数量的变化：
\[
\frac{dB}{dt}=rF\left(1-\frac{B}{K}\right)-\bar B
\]
刚刚进入成体阶段的七鳃鳗将经历性别分化。我们用 $\bar F$ 与 $\bar M$ 分别表示分化为雌性与雄性的个体数，它们显然与 $\bar B$ 满足如下关系：
\begin{equation}
\bar B=\bar F+\bar M \tag{1}
\end{equation}
于是雌性与雄性七鳃鳗的增长速率可描述为
\[
\begin{cases}
\dfrac{dF}{dt}=\bar F-\mu F\\[0.5em]
\dfrac{dM}{dt}=\bar M-\mu M
\end{cases}
\]
其中 $\mu$ 表示成年七鳃鳗的死亡率。

\subsection{差分方程形式}

七鳃鳗的种群动态具有如下特征：

\begin{itemize}
\item 环境变化具有周期性。与其他生物所处的环境一样，温度、pH 值等环境因素的变化以 12 个月为周期。
\item 七鳃鳗自身变化具有周期性节律。根据相关网站[6]，七鳃鳗的繁殖、孵化等行为必须在一年之中的特定时段内进行。
\end{itemize}

基于上述特征，将 3.1 节的微分方程转换为差分方程，可以更好地描述七鳃鳗的种群动态。相应地，我们采用
\begin{equation}
B_{t+1}-B_t=rF_t\left(1-\frac{B_t}{K}\right)-\bar B_t \tag{2}
\end{equation}
描述幼体期，采用
\begin{equation}
F_{t+1}-F_t=\bar F_t-\mu F_t \tag{3}
\end{equation}
描述雌性成体期，采用
\begin{equation}
M_{t+1}-M_t=\bar M_t-\mu M_t \tag{4}
\end{equation}
描述雄性成体期。所有差分方程的步长均为一个月。

增长率 $r$ 与环境容纳量 $K$ 都是影响种群动态的参数，而准确掌握七鳃鳗度过幼体期所需的时间同样至关重要，唯其如此才能确定第 $t$ 个月完成变态的个体数，即 $\bar B_t$。因此我们用 $\tau_i$ 表示第 $i$ 个月出生的七鳃鳗成长为成体七鳃鳗所需的时间。

在第 $t$ 个月成为成体的七鳃鳗，其出生月份 $i$ 与发育时长 $\tau_i$ 必须满足关系 $t=i+\tau_i$。于是，把所有满足该关系的七鳃鳗数量相加，即可得到第 $t$ 个月成熟的七鳃鳗数量，即：
\begin{equation}
\bar B_t=\sum_{\substack{t=i+\tau_i\\ 1\leqslant i\leqslant t}} rF_{i-1}\left(1-\frac{B_{i-1}}{K}\right) \tag{5}
\end{equation}
随后它们立即经历性别分化阶段，成为雌性或雄性。我们用 $\alpha_t$ 定义第 $t$ 个月的性别分化比例，即 $\alpha_t=\bar F_t/\bar B_t$。结合式 (1)，可得如下方程：
\begin{equation}
\bar F_t=\sum_{t=i+\tau_i}\alpha_i rF_{i-1}\left(1-\frac{B_{i-1}}{K}\right) \tag{6}
\end{equation}
\begin{equation}
\bar M_t=\sum_{t=i+\tau_i}(1-\alpha_i)\,rF_{i-1}\left(1-\frac{B_{i-1}}{K}\right) \tag{7}
\end{equation}
根据假设 5，性别分化比例 $\alpha_t$ 可由发育所消耗的时间决定，因此 $\alpha_t$ 与 $\tau_t$ 之间的关系可表示为
\[
\alpha_t=f(\tau_t).
\]
在参数恒定且取值适当的情况下，BFM 模型的仿真结果如图 4 所示：其中图 4a 对应 $\tau$ 与 $\alpha$ 均取常数的情形，图 4b 对应二者随时间变化的情形，模型均能给出结构完整的种群演化轨迹。

\section{任务 1 与任务 4：七鳃鳗与其他物种}

\subsection{$\alpha$ 与 $\tau$ 的确定}

根据文献[7]，$\alpha$ 仅与 $\tau$ 有关，且二者满足线性关系，即：
\[
\alpha_t=c\tau_t+d.
\]
通过数据拟合，我们得到参数取值 $c=-0.009$ 与 $d=0.764$。

接下来讨论 $\tau_t$ 的表达式。

对幼体七鳃鳗而言，成长为成体所需的平均能量可以近似认为彼此相等。我们假定，一旦某条幼体七鳃鳗累积到一定能量 $W$，它便能够成长为成体七鳃鳗。因此，七鳃鳗的固定能量等于摄入能量与未吸收能量之差。当累积的固定能量达到某一数值 $W$ 时，第 $t$ 个月出生的幼体七鳃鳗即完成生长，其累积时长即为 $\tau_t$。对单条幼体七鳃鳗而言，这一关系可表示为：
\[
\sum_{t+1\leqslant j\leqslant t+\tau_G}\left(\frac{food}{B_j}-waste\right)\geqslant W,\qquad \tau_t=\min\tau_G
\]
其中 $food$ 表示每月的食物供给量，$waste$ 表示未被吸收的能量。我们假定吸收率 $\eta$（即食物能量的同化效率）为常数，这意味着全部幼体七鳃鳗每月固定的能量为 $\eta\cdot food$。于是上述不等式可改写为：
\[
\sum_{t+1\leqslant j\leqslant t+\tau_G}\frac{1}{B_j}\geqslant\frac{W}{\eta\cdot food}
\]
当仿真月份 $t$ 达到一定规模（例如 $t>100$）时，幼体七鳃鳗的数量趋于稳定。因此我们可以近似认为 $B_t$ 不随时间变化。在此情形下，上述不等式可再次改写为：
\[
\tau_G\geqslant\frac{W}{\eta\cdot food\cdot B_t}
\]
考虑到当 $B_t\rightarrow 0$ 时会出现 $\tau_t\rightarrow 0$ 这一不合理结果，我们加入一个截距，使 $\tau_t$ 与 $B_t$ 满足线性关系：
\[
\tau_t=kB_t+b
\]
其参数取值为 $k=-0.000844$ 与 $b=35.678$。

\subsection{种间关系}

在特定生态系统中，七鳃鳗与其他物种以及非生物环境相互作用。例如，七鳃鳗作为入侵物种已对五大湖生态系统产生了显著影响。下面我们将分别讨论七鳃鳗作为猎物、寄生物与宿主时的数学模型，并简要分析其与非生物环境之间的关系。

\subsubsection{与天敌的关系}

首先讨论七鳃鳗与其天敌之间的模型。我们只引入一种天敌，以白斑狗鱼（Esox lucius）为例。

基于洛特卡–沃尔泰拉模型——该模型描述两个物种相互作用（一为捕食者、一为其猎物）的生态系统动态[8]——我们建立了按性别分离的改进型洛特卡–沃尔泰拉模型。模型的表达式如下：
\begin{align}
F_{t+1}-F_t&=\bar F_t-\mu F_t-\varepsilon_F F_t P_t \tag{8}\\
M_{t+1}-M_t&=\bar M_t-\mu M_t-\varepsilon_M M_t P_t \tag{9}\\
P_{t+1}-P_t&=-\mu_P P_t+\theta_F F_t P_t+\theta_M M_t P_t \tag{10}
\end{align}
其中：

\begin{itemize}
\item $P_t$ 表示第 $t$ 个月白斑狗鱼的种群数量。
\item $\mu_t$ 表示第 $t$ 个月白斑狗鱼的天然死亡率。
\item $\theta$ 描述七鳃鳗的单位能量；$\theta$ 取值越大，表示单位七鳃鳗所含能量越多。
\item $\varepsilon$ 描述被捕食的难度；$\varepsilon$ 取值越大，表示七鳃鳗越容易被捕食。
\end{itemize}

最为重要的是，我们对七鳃鳗区分了性别，因为我们假定雌性与雄性在单位能量与被捕食难度上存在性别差异。

图 6 展示了七鳃鳗与作为捕食者的白斑狗鱼的种群动态。由此可以得知：

\begin{itemize}
\item 捕食者与七鳃鳗的种群数量呈周期性波动，且捕食者种群的变化滞后于七鳃鳗。
\item 与 $\alpha$ 保持常数、种群在 $t=200$ 附近灭绝的情形相比，当 $\alpha$ 能够适应性变化时，七鳃鳗可以在一定范围内稳定存活更长时间。
\end{itemize}

\subsubsection{与宿主的关系}

其次，我们探讨七鳃鳗与其宿主之间的关系。本节中七鳃鳗作为寄生物，我们选取的宿主为鲫鱼。

我们采用尼科尔森–贝利模型刻画鲫鱼的种群动态[9]。不过，我们用自行建立的 BFM 模型而非尼科尔森–贝利模型来刻画（七鳃鳗的）种群动态，因为在现实情形中鲫鱼并不是唯一的宿主。

描述鲫鱼种群动态的表达式如下：
\begin{equation}
H_{t+1}=r_H H_t e^{-(\beta_F\cdot F_t+\beta_M\cdot M_t)} \tag{11}
\end{equation}
其中：

\begin{itemize}
\item $H_t$ 表示第 $t$ 个月鲫鱼的种群数量。
\item $r_H$ 表示鲫鱼单独存在时的增长率。
\item $\beta$ 刻画七鳃鳗对鲫鱼的影响程度。在尼科尔森–贝利模型中，$\beta$ 越大表明鲫鱼受损越严重。出于与 4.2.1 节几乎相同的理由，我们对七鳃鳗区分性别，因此雌性与雄性七鳃鳗造成的损害程度不同，分别记为 $\beta_F$ 与 $\beta_M$。
\end{itemize}

图 7 展示了七鳃鳗与作为宿主的鲫鱼的种群动态。由此可以得知：

\begin{itemize}
\item 当 $\alpha$ 保持常数时，宿主在较短时间内即被寄生至灭绝，这不利于七鳃鳗自身的长久存活；当 $\alpha$ 作为变量时，该情形可得到显著改善。
\item 当 $\alpha$ 为变量时，七鳃鳗能够与其宿主共存。这意味着二者之间可能发生协同演化。
\end{itemize}

\subsubsection{与寄生虫的关系}

反过来，现在轮到七鳃鳗被寄生了。我们选取一个体型较小的类群——疮痂鱼虱属（Lepeophtheirus spp.）来寄生七鳃鳗。

根据文献[10]，我们作如下假设：

\begin{itemize}
\item 不考虑幼体七鳃鳗被寄生的情形。
\item 雌性七鳃鳗更易被寄生。某些寄生虫只侵扰雌性七鳃鳗，因为它们偏好栖息于卵巢等部位；同时一般认为雌性具有更丰富的营养价值。
\end{itemize}

在此情形下，我们使用尼科尔森–贝利模型的另一侧，即描述疮痂鱼虱属种群动态的部分：
\begin{equation}
J_{t+1}=F_t\left(1-e^{-\gamma_F\cdot J_t}\right)+M_t\left(1-e^{-\gamma_M\cdot J_t}\right) \tag{12}
\end{equation}
其中：

\begin{itemize}
\item $J_t$ 表示第 $t$ 个月疮痂鱼虱属的种群数量。
\item $\gamma$ 刻画寄生虫对七鳃鳗的影响程度。在尼科尔森–贝利模型中，$\gamma$ 越大表明疮痂鱼虱属的寄生能力越强。同样地，由于雌性与雄性七鳃鳗的抵抗力存在差异，我们对七鳃鳗区分性别，并分别用 $\gamma_F$ 与 $\gamma_M$ 刻画其抵抗力。
\end{itemize}

图 8 展示了七鳃鳗与其寄生虫的种群动态。由此可以得知：

\begin{itemize}
\item 寄生虫的数量随七鳃鳗数量的变化而变化，寄生虫数量的变化略有滞后。
\item 当 $\alpha$ 固定时，寄生虫因数量过低而灭绝；而当 $\alpha$ 为变量时，寄生虫与七鳃鳗种群均保持稳定，能够长期共存。
\end{itemize}

\subsection{与非生物环境的关系}

除生态系统中的生物组分之外，七鳃鳗对生态系统内的非生物环境也有一定影响。幼体七鳃鳗会钻入沙质底质，从而影响底部的土壤状况。人类为控制七鳃鳗幼体而在湖底铺设砾石等材料的措施，进一步加剧了对非生物环境的影响。除此之外，七鳃鳗的其他发育阶段也会影响生态系统的环境。

\subsection{物种多样性指数}

在分析了七鳃鳗对生态系统中不同角色的影响之后，我们转而评价七鳃鳗所在生态系统中物种多样性的变化。按照相关文献，辛普森多样性指数与香农–维纳指数是评价生物多样性的两个良好指标。辛普森多样性指数表示从群落中随机抽取两个个体、它们属于同一物种的概率；香农–维纳指数则基于物理学与信息论中的熵概念来描述生态系统的物种多样性。二者的定义如下：
\begin{equation}
D=1-\sum_{j=1}^{Species}p_j^2 \tag{13}
\end{equation}
\begin{equation}
H=-\sum_{j=1}^{Species}p_j\cdot\log_2 p_j \tag{14}
\end{equation}
我们分别计算了性别比例恒定与性别比例可变的七鳃鳗种群所在生态系统的物种多样性。我们选取辛普森指数，其随时间变化的趋势如下：

\begin{itemize}
\item 在图 9a 中，辛普森指数在约 200 个月处均值骤降、振幅增大，变得不稳定。这是由于在 4.2.3 节的仿真结果中，寄生性七鳃鳗的寄生虫在一段时间后灭绝，导致辛普森指数的变化趋势在该时刻发生突变。
\item 在图 9b 中，辛普森指数在 0.65 附近波动，波形相对稳定、振幅较小，表明该种群所在生态系统的物种多样性能够维持在更高、更稳定的水平。
\end{itemize}

比较两图可以得出结论：具有可变性别比例的七鳃鳗种群能够提升生态系统的物种多样性。这是七鳃鳗带给生态系统中其他物种的一项优势。在 5.3.2 节中，我们将进一步利用物种多样性指数计算生态系统的稳定性。
\section{任务 2 与任务 3：生态系统中的七鳃鳗}

在上一节中，我们讨论了具有可变性别比例的七鳃鳗种群对生态系统中其他成员的影响。本节我们把关注点转向七鳃鳗种群自身，分析具有可变性别分化比例在生存上的优势与劣势。

\subsection{优点与缺点}

\subsubsection{适宜环境下}

首先，我们在环境适宜、资源充足、既无竞争也无捕食者的条件下模拟了两个七鳃鳗种群。对其中一个种群，性别分化比例 $\alpha$ 取为不随时间变化的常数，我们将其设为 0.45；对另一个种群，$\alpha$ 随环境变化而变化。结果如图 10 所示。

从图 10 的结果可以得知：

\begin{itemize}
\item 在第一种情形中，幼体七鳃鳗的数量迅速增长至环境容纳量，随后围绕均值波动，且振幅不断增大；成年七鳃鳗的数量同样围绕其均值波动且振幅不断增大，直至波动幅度过大，导致幼体数量降为零，七鳃鳗种群灭绝。这表明，即便在不存在竞争、资源充足的环境中，性别比例固定为 0.45 的七鳃鳗种群在生存上也是极不稳定的。
\item 在第二种情形中，幼体以及成年雄性、成年雌性七鳃鳗的数量在经历短暂增长后趋于稳定，围绕各自均值维持规律而持续的波动。这表明，具有可变性别比例的七鳃鳗种群具有更强的稳定性与长期存活的潜力。事实上，现实七鳃鳗种群的性别比例能够随环境适应性变化，从而增强其生存的稳定性。我们的数值仿真结果充分说明了这一现象。
\item 优点：稳定性更强。两种情形的比较表明，在资源充足的环境中，具有可变性别分化比例的七鳃鳗种群表现出更强的稳定性。
\end{itemize}

\subsubsection{扰动环境下}

接下来，我们假定环境的剧烈变化主要作用于死亡率。随后我们模拟了两个七鳃鳗种群对剧烈扰动的响应。我们将随时间变化的死亡率建模如下：
\begin{equation}
\mu_t=A_0+A_1\cdot\sin\left(\frac{2\pi}{T}\right)+A_2\cdot\delta\left(t-t_0\right) \tag{15}
\end{equation}
其中：

\begin{itemize}
\item $A_0$ 表示基线死亡率。
\item $A_1$ 用于调整由温度、繁殖周期等具有年周期性的因素所引起的死亡率的季节变化。
\item $A_2$ 用于刻画由环境突变或人类活动（例如集中捕捞七鳃鳗）所引起的死亡率骤增；$t_0$ 表示这类突发事件发生的时刻。
\end{itemize}

我们通过令 $t_0$ 取随机值来设置随机扰动。由于相关方程此前已经离散化，函数形式的参数化实现起来十分直接。结果如图 11 所示。

从图 11 可以得出：

\begin{itemize}
\item 在图 11a 中，遭受扰动后七鳃鳗种群数量显著下降，但很快恢复到扰动前水平，表明性别分化比例恒定的七鳃鳗种群具有一定的抗扰动能力。
\item 在图 11b 中，种群在扰动后同样显著下降并迅速恢复，但恢复后的种群动态与扰动前状态略有差异。这一现象说明，具有可变性别分化比例的七鳃鳗种群对扰动同样具备恢复力。
\item 具有可变性别分化比例的种群的优势在于其偏离正常状态的程度更小、恢复更快，表明其在扰动下具有更强的稳定性。为食用目的对七鳃鳗进行小规模捕捞，能够使七鳃鳗种群得以持续存活。本模型将在 5.3 节中对这种抗扰动能力作进一步细化与定量分析。
\item 当人类干预旨在控制入侵性七鳃鳗种群以维持生态平衡时，这种恢复力可能构成劣势：其稳定性使它们难以被治理，进而可能对当地渔业与湖泊经济产生负面影响。
\end{itemize}

\subsubsection{繁殖成功率}

具有可变性别分化比例的七鳃鳗种群可能导致最终的性别比例（即雌性与雄性的实际比值）偏离 0.5。按照费希尔性别比理论（Fisher's sex ratio theory），在自然界中经过充分演化，雌雄比例应当趋于 1:1。显著偏离该比例会导致数量较多的性别一方在择偶上的竞争加剧，交配机会减少，不利于基因传递。这可能使携带有利基因突变的个体因错失交配机会而无法将其传递下去[11]。个体获得交配机会的可能性既取决于自身性别的比例，也取决于异性的比例。因此，繁殖成功率定义如下：
\begin{equation}
G=Rate_F\cdot Rate_M \tag{16}
\end{equation}
我们计算了不同时刻的 $G$ 值，从而得到一系列随时间变化的繁殖成功率。对这些值取平均后，结果如下：当 $\alpha$ 为常数时，$G=0.2467$；当 $\alpha$ 为变量时，$G=0.2114$。

\begin{itemize}
\item 劣势：繁殖成功率更低。具有可变性别比例的七鳃鳗种群面临更低的繁殖成功率。这表明，即便在最适宜的环境中，仍有相当数量的个体可能得不到繁殖机会，其基因无法传递下去。从长期看，这可能对该物种的演化不利。七鳃鳗是现存最古老的无颌脊椎动物之一，其历史可追溯至奥陶纪；它们随时间的推移演化速率放缓，保留了使其被称为“活化石”的原始特征。这一发现与我们的分析相符。
\end{itemize}

\subsection{反馈模型}

基于 5.1 节的分析，具备调整性别比例能力的七鳃鳗种群在稳定环境与扰动环境中都表现出更高的稳定性。为探究其原因，下一步需对方程进行分析，以解释可变性别比例所带来的稳定性。因此我们建立反馈模型，对这一现象作进一步讨论。

下面我们对具有可变性别比例的七鳃鳗种群作一简要分析：假定 BFM 模型的各成员在某一时刻处于均衡状态，此时幼体数量 $B$ 的扰动会提高人均营养，加快发育并缩短发育时间。这会抬高性别比例，使本代幼体中分化为雌性的数量增加，从而增强七鳃鳗种群的繁殖能力。下一代幼体补偿了最初的减少，通过负反馈过程使幼体种群在外部扰动下保持稳定。

为进一步阐明可变性别比例七鳃鳗种群对外部扰动的抵抗力，我们为两类七鳃鳗种群分别引入反馈模型。我们在图 12 中直观地表示了性别比例恒定与可变两种情形下，幼体数量、成年雌性与雄性数量以及性别分化比例 $\alpha$ 之间的关系。

在图中，乘法器的输出为其两路输入之积。我们令 $\alpha$ 保持为常数 0.45。这会使雌性比例随时间推移而下降，从而加剧原有的趋势。这一过程在控制系统中构成正反馈。在现实中，正反馈并不总是导致种群数量持续下降；相反，在适当的输入下（由本模型中 $F$ 与 $M$ 的初值所代表），它会产生振幅逐渐增大的自激振荡。在图 10a 中，$\alpha$ 为常数，种群呈现出振幅随时间不断增大的波动。图 10a 的结果与我们的正反馈分析一致，进一步验证了该分析的有效性。

图 12b 与图 12a 的差别在于：在图 12b 中，幼体种群能够调整 $\alpha$ 的取值，从而形成一个新的反馈回路，构成负反馈过程。根据控制理论，负反馈系统能够在一定范围内抵消外部扰动，使系统的输出趋于稳定。如图 11b 所示，七鳃鳗种群在遭受剧烈环境扰动后能够存活并很快恢复到原有水平，这与我们的负反馈分析相符。

我们将参数 $\alpha$ 设为 0.85，进行数值仿真，随后在图 13 中绘制 $F-B-t$ 相轨迹随时间的变化，以观察其演变。

从图 13 可以得出以下结论：

\begin{itemize}
\item 在图 13a 中，相轨迹的半径逐年呈增大趋势，形成“喇叭”形。这表明七鳃鳗的种群振荡逐年加剧、不稳定程度不断恶化，与我们的正反馈模型分析一致。
\item 在图 13b 中，经过短暂的暂态过程后，相轨迹的半径逐渐收敛到一个稳定的“极限环”，其形状与振幅均保持稳定，与我们的负反馈模型分析一致。
\item 优点：稳定性得到强化。具有可变性别比例的七鳃鳗种群的负反馈特性，增强了其在正常环境以及面对扰动时的稳定性。
\item 劣势：在适宜条件下难以繁荣。同样的负反馈特性可能限制七鳃鳗种群适应新的、更有利环境的能力，使其维持现状，并可能抑制其适应能力。
\end{itemize}

\subsection{对生态系统稳定性的影响}

为分析七鳃鳗种群对其所在生态系统稳定性的影响，我们首先需要对生态系统稳定性作出定量描述。生态系统稳定性可以从多个角度加以描述。

\subsubsection{抵抗力稳定性与恢复力稳定性}

生态系统稳定性可分为两部分：抵抗力稳定性与恢复力稳定性[12]。

\begin{itemize}
\item 抵抗力稳定性描述生态系统抵御外部扰动并维持原有状态的能力。
\item 恢复力稳定性描述生态系统在受到扰动后恢复到初始状态的能力。
\end{itemize}

为结合我们的数值仿真结果对这两类稳定性作定量描述，我们定义抵抗力稳定性 $S_1$ 为扰动下种群曲线与未扰动种群曲线之间的最大差值；恢复力稳定性 $S_2$ 定义为二者差值回到零所需的时间。具体定义如下：
\begin{equation}
S_1=D_\infty\left(f_1,f_2\right)=\sup_{t_1\leqslant t\leqslant t_2}\left|f_1(t)-f_2(t)\right| \tag{17}
\end{equation}
\begin{equation}
S_2=t_2-t_1 \tag{18}
\end{equation}
此后，我们引入综合稳定性 $S$ 的定义，即两条曲线之间所围成的面积。$S$ 的数学定义为：
\begin{equation}
S=\int_{t_1}^{t_2}\left|f_1(t)-f_2(t)\right|dt \tag{19}
\end{equation}
综合稳定性的概念同时涵盖抵抗力稳定性与恢复力稳定性。它们的取值越小，生态系统就越稳定。

利用 5.1.2 节的结果，我们计算其与正常状态之间的差异。结果如图 14 与图 15 所示。

由图 14 与图 15 可以清楚地看到：

\begin{itemize}
\item 在加入扰动之前，两条曲线完全重合；加入扰动之后，两条曲线呈现不同的走势，但经过一段时间后，受扰动的曲线重新回到正常水平。
\end{itemize}

接下来，我们分别计算两个七鳃鳗种群在两次扰动下的稳定性，结果如下：

\begin{itemize}
\item 当 $\alpha$ 为常数时：第一次扰动的 $S_1=1175.7$、$S_2=90$、$S=51070$；第二次扰动的 $S_1=1162$、$S_2=100$、$S=53362$；平均值为 $S_1=1168.8$、$S_2=95$、$S=52216$。
\item 当 $\alpha$ 为变量时：第一次扰动的 $S_1=1238.4$、$S_2=83$、$S=61324$；第二次扰动的 $S_1=913.612$、$S_2=75$、$S=34078$；平均值为 $S_1=1076$、$S_2=79$、$S=47701$。
\end{itemize}

以两次扰动的平均稳定性作为评价标准，并结合上述细节，可以得出：

\begin{itemize}
\item 具有可变性别比例的七鳃鳗种群表现出更强的抵抗力稳定性、恢复稳定性与整体稳定性。这表明具有变化性别比例的七鳃鳗种群能够增强生态系统的稳定性。
\end{itemize}

\subsubsection{生物多样性}

生物多样性可以间接反映生态系统稳定性的大小。我们引入辛普森多样性指数与香农–维纳指数，因为二者都是评价生物多样性的良好指标。我们计算了七鳃鳗、捕食者、寄生虫与宿主之间相互作用所形成的群落的这两个指标值，结果如下：当 $\alpha$ 为常数时，$H_{mean}=1.4010$、$D_{mean}=0.4978$；当 $\alpha$ 为变量时，$H_{mean}=1.9655$、$D_{mean}=0.6752$。

上述结果表明，具有可变性别比例的七鳃鳗种群提高了其所在环境的生物多样性。这说明七鳃鳗能够增强生态系统的稳定性。

\section{灵敏度分析}

在我们的模型中，还有一些与真实情形密切相关的参数。在此前的计算中，这些参数的取值来自特定情形，但其数值可能随情形的变化而变化。我们可以改变这些参数的取值，以计算模型对这些参数的灵敏度。通过对这些参数进行灵敏度分析，我们得到如下结论。

\subsection{七鳃鳗宿主的自然增长率 $r_H$}

$r_H$ 表示在七鳃鳗未入侵时七鳃鳗宿主的自然增长率。七鳃鳗的宿主包括许多不同种类的鱼，各自的 $r_H$ 取值不同。影响 $r_H$ 的因素包括出生率与死亡率；$r_H$ 越高，表明该鱼类的繁殖能力与环境适应能力越强。

在 4.2.2 节中，我们讨论了七鳃鳗的数量与性别变化对生态系统中宿主种群数量的影响。研究发现，当七鳃鳗的性别分化率可以变化时，宿主与七鳃鳗之间会建立起相对稳定的数量关系。然而，当七鳃鳗入侵某一生态系统时，其与宿主的关系会变得更加复杂。

\begin{itemize}
\item 在图 16a 中（$r_H=1.0350$），宿主种群受到七鳃鳗入侵的巨大影响，在该生态系统中，该物种于入侵后约第 400 个月发生不可逆的灭绝。这类宿主的繁殖能力与环境适应能力较弱，无法有效抵御外来物种入侵对本地生态环境造成的冲击。
\item 在图 16b 中（$r_H=1.0465$），宿主种群在七鳃鳗入侵后迅速减少，但经过一段时间的恢复后稳定在远低于入侵前的水平。这类宿主受七鳃鳗入侵影响较大，但经过一段时间后适应了七鳃鳗的寄生，达到相对稳定的种群数量。然而，入侵造成的影响无法完全恢复，种群数量只能维持在较低水平。
\item 在图 16c 中（$r_H=1.0475$），宿主种群在七鳃鳗入侵初期略有减少，但经过一段时间后维持在高于入侵前的水平。这类宿主受七鳃鳗入侵的影响较小，在与七鳃鳗长期共存的过程中维持了较高的种群数量。我们推测，这是由于宿主与七鳃鳗在长期中发生了协同演化，同时该宿主的其他竞争者受七鳃鳗的影响更大、竞争力被削弱，从而更有利于该宿主种群的扩张。
\end{itemize}

这表明宿主种群数量对 $r_H$ 的变化高度敏感。这说明七鳃鳗作为入侵物种，其寄生行为对本地物种及原生生态系统具有显著影响。因此，必须寻求防控七鳃鳗的策略。

\subsection{幼体七鳃鳗的环境容纳量 $K$}

幼体七鳃鳗的容纳量 $K$ 与其可获得的空间与环境资源总量有关，也与其生存环境的温度、含氧量和 pH 值是否适宜有关。通过对 $K$ 进行灵敏度分析，我们可以判定幼体七鳃鳗对不同环境的适应能力。

由图 17 的数据可以得出结论：当 $K$ 变化时，我们的 BFM（生物能量鱼类模型，Bioenergetic Fish Model）的增长趋势几乎保持不变，对 $K$ 的变化表现出较低的灵敏度与良好的稳定性，这是我们模型的一个优点。三图的比较表明，七鳃鳗的最大种群规模受 $K$ 变化的影响显著，但均值大致相同；$K$ 越大，种群稳定性越低。由于人工剔除幼体七鳃鳗十分困难，生长控制措施通常包括阻止幼体七鳃鳗钻入底质，以及使用 TFM（杀七鳃鳗剂）通过降低 $K$ 来有效治理七鳃鳗入侵[1]。

接下来我们将讨论如何控制成年七鳃鳗。

\subsection{成年七鳃鳗的死亡率 $\mu$}

在 5.1.2 节中，我们以式 (15) 提出了成年七鳃鳗的死亡模型。通过对 $A_0$、$A_1$ 与 $A_2$ 的灵敏度分析，可以得出结论：七鳃鳗种群对死亡率的变化高度敏感。

\begin{itemize}
\item 在图 10b 中，参数取值为 $A_0=0.05$、$A_1=0.05$、$A_2=0$，其中 $A_0$ 与 $A_1$ 相对较低，可视为在没有任何导致大规模死亡的灾难性事件时成年七鳃鳗的自然死亡率。
\item 在图 18a 中，$A_0=0.15$，$A_1$ 与 $A_2$ 保持不变，表示自入侵之初起通过有计划的持续捕捞来控制七鳃鳗入侵。此时七鳃鳗种群的不稳定性显著增强，雌鱼数量在第 38 年前后过低，导致种群灭绝。
\item 在图 18b 中，$A_1=0.25$，$A_0$ 与 $A_2$ 保持不变，表示自七鳃鳗入侵之初起按年周期进行生物防治或化学处理。由于自然状态下七鳃鳗数量本就存在一定周期的波动，季节性死亡率的上升使种群数量随时间的变化更加剧烈，进一步加剧了种群的不稳定性与灭绝的可能性。图中雌鱼数量在第 33 年降至最低点，因数量不足而导致种群灭绝。
\item 在图 18c 中，$A_2=0.3$，$A_0$ 与 $A_1$ 保持不变，表示在某一年对成年七鳃鳗进行大规模捕捞。通过改变 $A_2$ 发生变化的时间点，我们观察到：在成年雌性七鳃鳗数量处于最低值时集中进行大规模捕捞效果最佳。图 18c 中的集中捕捞发生在第 53 年，此时成年雌性七鳃鳗数量处于最低水平，约捕获 90\% 的成体。然而，由于幼体无法被捕捞，种群仍可继续繁殖。约在第 72 年，成年七鳃鳗数量出现骤增，雄性数量远高于捕捞前水平。这是因为单次过度捕捞导致七鳃鳗幼体数量减少，改变了发育时长 $\tau$，使部分幼体在同一时段完成分化，从而引起数量反弹。
\end{itemize}

综上，七鳃鳗种群对死亡率的变化高度敏感。控制七鳃鳗需要同时针对幼体与成体。

\section{模型评价}

\subsection{优点}

\begin{itemize}
\item 我们细致地刻画了七鳃鳗种群的性别比例随环境条件发生适应性变化的内在机制。
\item 我们从方程本身出发，结合负反馈机制，建立了能够解释具有可变性别比例的七鳃鳗种群为何具备更强适应性的反馈框架。
\item 我们全面考虑了七鳃鳗种群在不同情形下的生存状况，并选取了恰当的指标进行定量表征。
\item 我们的灵敏度分析结果既可用于讨论模型自身的稳定性，也可为七鳃鳗生物入侵的防控提供具有实际意义的指导。
\end{itemize}

\subsection{缺点与进一步讨论}

\begin{itemize}
\item 我们建立的反馈模型是一个半定量模型。尽管它能够很好地解释各项任务，但其中包含了定性描述，因而无法精确判断系统所处反馈的类型。系统是否构成反馈取决于反馈系统的具体参数，例如反馈回路的反馈系数。事实上，幼体七鳃鳗数量 $B$ 与环境容纳量 $K$ 也可以构成一个负反馈系统。这些相互嵌套的控制回路所组成的系统，远比我们的定性分析复杂。
\item 我们的模型可以从这一角度加以改进：对方程所描述的微分动力系统开展定量研究，并结合相关理论，计算参数 $r$、$K$ 的取值与种群数量是否处于反馈状态之间的关系。
\end{itemize}

\section*{参考文献}

\noindent[1] Great Lakes Fishery Commission, Sea lamprey -- an ancient invader, http://www.glfc.org/sea-lamprey.php, Accessed: February 5, 2024, 2023.

\noindent[2] Y. Wang, ``Stochastic evolution analysis of asymmetric game models,'' M.S. thesis, 2024.

\noindent[3] N. S. Johnson, T. J. Buchinger, and W. Li, ``Reproductive ecology of lampreys,'' Lampreys: Biology, Conservation and Control: Volume 1, pp. 265--303, 2015.

\noindent[4] G. J. Farmer, ``Biology and physiology of feeding in adult lampreys,'' Proc. Biol. Sci., vol. 284, no. 1851, p. 20 170 262, Mar. 2017.

\noindent[5] S. Xue, ``Study on the population control and evolution using sex- and stage- structured mathematical models,'' Ph.D. dissertation, Donghua University, 2023.

\noindent[6] Taonga species: Piharau, https://niwa.co.nz/TaongaSpecies\_PiharauLOWRES.pdf.

\noindent[7] N. S. Johnson, W. D. Swink, and T. O. Brenden, ``Field study suggests that sex determination in sea lamprey is directly influenced by larval growth rate,'' Proc. Biol. Sci., vol. 284, no. 1851, p. 20 170 262, Mar. 2017.

\noindent[8] S. E. J\o{}rgensen and B. D. Fath, Encyclopedia of Ecology. 2008.

\noindent[9] R. E. Ricklefs, The Economy of Nature, 5th ed. W.H. Freeman, 2001.

\noindent[10] M. A. Shavalier, M. Faisal, M. L. Moser, and T. P. Loch, ``Parasites and microbial infections of lamprey (order petromyzontiformes berg 1940): A review of existing knowledge and recent studies,'' Journal of Great Lakes Research,

\noindent[11] C. Niu, A. Lou, R. Sun, and Q. Li, Foundations in Ecology. 2015, ISBN: 9787040423785.

\noindent[12] R. Sun, D. Wang, C. Niu, D. Liu, and L. Zhang, Principles of Animal Ecology, 4th ed. 2001.

\section*{人工智能使用报告}

1. OpenAI ChatGPT（2023 年 11 月 5 日版本，ChatGPT-4）

\noindent\textbf{提问 1：}我们发现用于描述生物多样性的香农–维纳指数与信息熵的表达式十分相似。请解释这两个概念之间的联系。

\noindent\textbf{输出：}常用于生态学研究中描述生物多样性的香农–维纳指数，与克劳德·香农（Claude Shannon）在信息论背景下提出的信息熵概念，确实关系密切，因为二者都在度量一个系统中的不确定性或多样性。二者的联系如下。

\noindent\textbf{信息熵}\quad 信息熵又称香农熵，它量化一组结果中所含不确定性或随机性的大小。信息熵 $H$ 的公式为：
\[
H=-\sum_{i=1}^{n}p\left(x_i\right)\log p\left(x_i\right)
\]
其中 $p(x_i)$ 是随机变量第 $i$ 个可能取值出现的概率，求和遍历该变量的全部可能取值。该式度量一条消息所包含信息的期望值（以比特为单位），其前提是假定所有可能的消息等概率出现。它实质上量化的是每传送一个信源符号平均所产生的信息量。

\noindent\textbf{生态学中的香农–维纳指数}\quad 香农–维纳指数（或称香农多样性指数）是香农熵概念在生物多样性上的应用。在生态学中，它用于量化一个群落内物种的多样性，同时考虑物种数目（丰富度）与其分布的均匀度。其公式与信息熵十分相似：
\[
H'=-\sum_{i=1}^{s}p_i\log p_i
\]
其中 $p_i$ 是数据集中属于第 $i$ 个物种的个体所占比例，$s$ 为物种总数。该指数实质上度量的是从数据集中随机抽取一个个体时，预测其所属物种的不确定性。香农–维纳指数越高，表明生态系统的多样性越高。

\noindent\textbf{联系}\quad 这两个概念的根本联系在于它们都着眼于度量某一分布中的不确定性、多样性或随机性。在信息论中，熵度量预测一条消息取值的不确定性；在生态学中，香农–维纳指数通过量化预测随机抽取个体所属物种的不确定性，来度量生态系统的生物多样性。二者都用对数函数对概率加权，这与信息或生物多样性随实体数目或物种数目及其分布而变化的缩放方式相符。这种数学上的相似性凸显了信息与生态多样性两个概念之间更深层的联系：无论是一组消息还是一个生物群落，二者都可被视为对一个系统中多样性程度与不可预测性的度量。

\end{document}
